\chapter{系统设计与实现}

% ===== 提示 =====
% 本章是论文的核心章节，详细描述系统的设计方案和实现过程。
% 根据你的实际项目调整内容，可以拆分为多个章节。
% ================

\section{系统需求分析}

\subsection{功能需求}

在此描述系统的功能需求，可以使用用例图等进行说明。

\subsection{非功能需求}

在此描述系统的性能、安全性、可用性等非功能需求。

\section{系统总体设计}

\subsection{系统架构设计}

在此描述系统的总体架构，可以插入架构图。插入图片的示例如下：

% 取消注释并替换为你的图片路径
% \begin{figure}[htbp]
%     \centering
%     \includegraphics[width=0.8\textwidth]{images/your_image.png}
%     \caption{系统架构图}
%     \label{fig:architecture}
% \end{figure}

% 引用图片：如图\ref{fig:architecture}所示。

\subsection{数据库设计}

在此描述数据库的设计方案，包括ER图和主要数据表结构。

\section{系统详细设计与实现}

\subsection{模块一实现}

在此描述各功能模块的具体实现方法和关键代码。

如需插入代码片段，可以使用\verb|verbatim|环境：

\begin{verbatim}
def hello_world():
    print("Hello, World!")
\end{verbatim}

\subsection{模块二实现}

在此继续描述其他功能模块的实现。

\section{系统测试}

\subsection{测试环境}

描述测试所使用的硬件和软件环境。

\subsection{功能测试}

描述功能测试的方案、用例和结果。

\subsection{测试结论}

总结测试结果，说明系统是否满足设计要求。
